% Reproducibility appendix. Every value below was read off the training code and the logged
% per-run configs (probe_config.json / seed-sweep JSONs). Items wrapped in \chk{...} disagree
% between the logs and tab_hyperparams.tex, or could not be verified -- resolve, then delete
% the macro.
\providecommand{\chk}[1]{\textcolor{red}{[CHECK: #1]}}

\subsection{Experimental setup and training details}
\label{app:hyperparam}

\paragraph{Shared protocol.}
All four pipelines are \emph{post-hoc}: the image backbone and the concept extractor are
frozen, concept activations $c(x)\in\mathbb{R}^{|\mathcal{C}|}$ are computed once and cached,
and only the final linear layer (baseline) or the gated head (GCG) is trained on the cached
activations. For each (model, dataset) pair, the baseline and the GCG head are trained on
\emph{identical} cached activations. Only the seed-dependent final stage (head initialisation
and minibatch order) is repeated for each seed $s\in\{0,10,20,30,40\}$. Before each run we
call \texttt{torch.manual\_seed}$(s)$, and for UCBM we also seed Python/NumPy/CUDA and set
cuDNN to deterministic mode. Seed-independent stages (concept discovery, the projection or
concept layer, activation caching) run once, using the seed of the original implementation.
We report top-1 accuracy on the official test split (Places-365: the official validation
split), as the mean $\pm$ std over the five seeds.

\paragraph{GCG head.}
The head is $\hat y = W\,(c(x)\odot m) + b$ with one gate logit $\mu_j$ per concept and
$g_j=\sigma(\mu_j/T)$, $T=1$. We use the \emph{hard-forward} straight-through estimator for
every run: the forward pass uses $m_j=\mathbb{1}[g_j>0.5]$ and the backward pass uses
$\partial m_j/\partial g_j=1$. All logits are initialised to $\mu_j=2$ ($g_j\approx0.88$, so
every gate starts open). The weights are initialised with the PyTorch \texttt{nn.Linear}
default. The per-minibatch objective is
\[
\mathcal{L} = \mathrm{CE}(\hat y, y) \;+\; \lambda_{\text{gate}}\sum_{j}\sigma(\mu_j)
\;+\; \mathcal{R}_{\text{model}}(W),
\]
where CE is averaged over the minibatch and the gate penalty is \emph{not} normalised by the
batch or by $|\mathcal{C}|$. $\mathcal{R}_{\text{model}}$ is the backbone-specific weight
penalty given below. It is zero for LF-CBM and VLG-CBM. We apply no warm-up of
$\lambda_{\text{gate}}$, no temperature annealing, and no early stopping; the head from the
last epoch is used. A concept is \emph{open} if $g_j>0.5$.

\paragraph{Mask-and-refit.}
After gated training we keep $K=\{j: g_j>\tau\}$ with $\tau=0.5$, delete every other
concept, and retrain a plain linear head on $c(x)_K$ with no gate and no regulariser. The head
is warm-started from the gated weights with the gate folded in ($W_{:,K}\odot g_K$). We use
Adam with learning rate $10^{-3}$, no weight decay, 20 epochs, the same batch size as gated
training, and refit seed $=s$. The refit head has exactly $|K|$ non-zero columns.

\paragraph{Counting concepts in the baselines.}
The sparse baselines have no gate, so a concept counts as used if any class has
$|W_{kj}|>10^{-5}$ on it (LF-CBM, UCBM, VLG-CBM), or if its weight clears the
\texttt{concept\_usage\_report} threshold for DN-CBM \chk{state the DN-CBM threshold}.

\subsubsection{Per-backbone details}

\paragraph{DN-CBM.}
\emph{Concepts:} We use the released CLIP-RN50 sparse autoencoder of DN-CBM (1024-d image
embedding $\to$ 8192 latents). Latents are named by their nearest CLIP text embedding in the
20k CLIP-Dissect vocabulary. \emph{Data:} From the training split we hold out a stratified 20\%
(split seed 42). It is used only to select $\lambda_{\text{gate}}$, and all heads are trained
on the remaining 80\%.
\emph{Baseline:} A linear probe with $\lambda_{\ell_1}\sum_{kj}|W_{kj}|$, $\lambda_{\ell_1}=10^{-4}$,
Adam at a constant learning rate of $10^{-3}$, 100 epochs.
\emph{GCG:} Same optimiser and learning rate, with $\mathcal{R}_{\text{model}}=\lambda_{\ell_1}\cdot
\mathrm{mean}_{kj}|W_{kj}|$ and $\lambda_{\ell_1}=10^{-4}$. Note that the gated arm uses the
\emph{mean}, not the sum, so its weight penalty is $|\mathcal{C}|\cdot K$ times weaker than
the baseline's. Epochs: 100 (CIFAR-100), 150 (CUB, Places-365).
\emph{Batch size:} 64 (CIFAR-100, CUB) and 4096 (Places-365, 1.44M training images).
\chk{The latest Places-365 baseline on disk used 150 epochs and $\lambda_{\ell_1}=10^{-5}$ (47.6\%);
the 100-epoch/$10^{-4}$ run gives 36.8\%. Say which one the paper reports.}

\paragraph{LF-CBM.}
\emph{Concepts:} GPT-3 concept sets from the original release (\texttt{<dataset>\_filtered.txt}).
Concept targets are CLIP ViT-B/16 image--text similarities. We drop concepts whose mean top-5
CLIP activation is $\le$ 0.25 / 0.26 / 0.28 (CIFAR-100 / CUB / Places-365), and concepts
whose projection similarity on the validation set is $\le 0.45$.
\emph{Projection:} A bias-free linear map from backbone features to concepts, trained to
maximise $\cos^3$ similarity to the CLIP targets. Adam, learning rate $10^{-3}$, batch
$\min(50000, n)$, at most 1000 steps, early-stopped on validation similarity (checked every 50
steps, patience 1). Concept activations are standardised with the training-set mean and std.
\emph{Backbones and layers:} CLIP-RN50 image encoder (CIFAR-100); CUB-trained ResNet-18,
\texttt{features.final\_pool} (CUB); Places-365 ResNet-50, \texttt{layer4} (Places-365).
Features are average-pooled.
\emph{Baseline:} GLM-SAGA elastic net ($\alpha=0.99$, step size 0.1, a single $\lambda$ on the
path), batch 64, with $\lambda$ / iterations of $7\times10^{-4}$/1000 (CIFAR-100),
$2\times10^{-4}$/5000 (CUB) and $3\times10^{-4}$/80 (Places-365). Weights and bias are
initialised to zero.
\emph{GCG:} Adam at a constant learning rate of $10^{-3}$, batch 64, no weight penalty.
Epochs: 100 (CIFAR-100), 150 (CUB), \chk{20 in the Places-365 hard-forward run on disk; the
caption says 150}.

\paragraph{UCBM.}
\emph{Concepts:} UCBM concept banks of 1000 / 200 / 3650 unsupervised concept vectors
(\texttt{concepts\_1000\_64}, \texttt{concepts\_200\_64}, \texttt{concepts\_3650\_64})
\chk{describe how the banks were extracted}. Concept activations are the cosine similarity of
backbone features to each concept vector, standardised when \texttt{normalize} is on
(CIFAR-100, CUB).
\emph{Backbones:} ImageNet ResNet-50 (torchvision V2 weights) for CIFAR-100, CUB ResNet-18,
Places-365 ResNet-18.
\emph{Classifier:} A per-concept learned scale $e^{s_j}$ ($s_j$ initialised to 0; disabled on
CUB) and offset $e^{o_j}$ ($o_j$ initialised to $-10$), then ReLU, then concept dropout ($p=0.1$
on CIFAR-100, $0.2$ otherwise), then a linear layer. Regularisers: an elastic net
($\alpha=0.99$) on post-ReLU activations with $\lambda_\pi=10^{-4}$, and an elastic net on $W$
with $\lambda_w=10^{-4}$. In the gated runs $\mathcal{R}_{\text{model}}$ keeps the $\lambda_w$
term, and the gate is applied after the ReLU and before dropout. Adam, learning rate
$10^{-3}$, batch 64.
\emph{Baseline:} Cosine learning-rate decay over 20 / 150 / 20 epochs.
\emph{GCG:} Constant learning rate. Epochs: 100 (CIFAR-100), 150 (CUB), 20 (Places-365). On
Places-365 the gated run sets $\lambda_\pi=0$.
\chk{The CUB recipe passes \texttt{--k 66}, but with the default ReLU selector the top-$k$
branch is never reached, so no top-$k$ is applied. Drop it from the description, or rerun with
\texttt{--relu no}.}

\paragraph{VLG-CBM.}
\emph{Concepts:} Concept sets from the original release, filtered by Grounding-DINO
annotations (confidence threshold 0.15). A concept-bottleneck layer (CBL) is trained with
multi-label BCE, Adam, learning rate $5\times10^{-4}$ and weight decay $10^{-5}$. Settings for
CIFAR-100 / CUB / Places-365: epochs 40/75/4, batch 32/32/128, hidden layers 1/0/1, positive
weight 0.2 (auto) / 1.0 / 50.0, crop-to-concept probability 0/0.5/0.5. The backbone is not
fine-tuned, and the shared-stage seeds are 28527 / 6885 / 42. CBL outputs are standardised with
training-set statistics. A fraction 0.1 / 0.1 / \chk{0.05 in config, 0.1 in the sweep recipe}
of the training set is held out for validation.
\emph{Backbones:} CLIP-RN50 (penultimate layer) for CIFAR-100, CUB ResNet-18
(\texttt{features.final\_pool}) for CUB, ResNet-50 (\texttt{layer4}) for Places-365
\chk{the table says CLIP RN50 for Places-365, but the config and logged args say ResNet-50}.
\emph{Baseline:} GLM-SAGA ($\alpha=0.99$, step size 0.1, batch 512) with
$\lambda=7\times10^{-4}$ / $2\times10^{-4}$ / $7\times10^{-4}$ and 2000 / 4000 / 200
iterations.
\emph{GCG:} Adam with learning rate $10^{-3}$, cosine-annealed over training, batch 512, no
weight penalty. The number of epochs equals the GLM-SAGA iteration count, i.e.\ 2000
(CIFAR-100) and 4000 (CUB) epochs.
\chk{No Places-365 VLG-CBM seed-sweep results are on disk; confirm $\lambda_{\text{gate}}$,
epochs, and that the run exists.}

\begin{table}[!h]
\centering
\small
\setlength{\tabcolsep}{3pt}
\caption{Full per-run settings. ``Base.\ reg.'' is the baseline's sparsity strength
(DN-CBM: $\ell_1$ sum; LF-CBM/VLG-CBM: GLM-SAGA elastic net; UCBM: $\lambda_w$/$\lambda_\pi$
elastic net). ``Base.\ it.'' is epochs (DN-CBM, UCBM) or GLM-SAGA iterations (LF-CBM,
VLG-CBM). All GCG heads use hard-forward STE, $\mu_j(0)=2$, $T=1$, Adam, and
mask-and-refit at $\tau=0.5$ (20 epochs, learning rate $10^{-3}$).}
\label{tab:hyperparams_full}
\begin{tabular}{llrlrrrrrl}
\toprule
& & & & \multicolumn{2}{c}{\textbf{Baseline}} & \multicolumn{4}{c}{\textbf{GCG}} \\
\cmidrule(lr){5-6}\cmidrule(lr){7-10}
\textbf{Model} & \textbf{Data} & $|\mathcal{C}|$ & \textbf{Backbone} & Base.\ reg. & Base.\ it. &
$\lambda_{\text{gate}}$ & Ep. & Batch & LR \\
\midrule
DN-CBM  & C100 & 8192 & CLIP RN50    & $10^{-4}$ & 100 & $5\!\times\!10^{-4}$ & 100 & 64   & $10^{-3}$ const \\
        & CUB  & 8192 & CLIP RN50    & $10^{-4}$ & 100 & $3\!\times\!10^{-3}$ & 150 & 64   & $10^{-3}$ const \\
        & P365 & 8192 & CLIP RN50    & \chk{} & \chk{} & $1\!\times\!10^{-4}$ & 150 & 4096 & $10^{-3}$ const \\
\midrule
LF-CBM  & C100 & 829  & CLIP RN50    & $7\!\times\!10^{-4}$ & 1000 & $1\!\times\!10^{-4}$ & 100 & 64 & $10^{-3}$ const \\
        & CUB  & 370  & ResNet-18    & $2\!\times\!10^{-4}$ & 5000 & $7\!\times\!10^{-4}$ & 150 & 64 & $10^{-3}$ const \\
        & P365 & \chk{2161} & ResNet-50 & $3\!\times\!10^{-4}$ & 80 & \chk{$1\!\times\!10^{-4}$} & \chk{20} & 64 & $10^{-3}$ const \\
\midrule
UCBM    & C100 & 1000 & ResNet-50    & $10^{-4}$/$10^{-4}$ & 20  & $1\!\times\!10^{-4}$ & 100 & 64 & $10^{-3}$ const \\
        & CUB  & 200  & ResNet-18    & $10^{-4}$/$10^{-4}$ & 150 & $5\!\times\!10^{-4}$ & 150 & 64 & $10^{-3}$ const \\
        & P365 & 3650 & ResNet-18    & $10^{-4}$/$10^{-4}$ & 20  & \chk{$1\!\times\!10^{-4}$} & 20 & 64 & $10^{-3}$ const \\
\midrule
VLG-CBM & C100 & 780  & CLIP RN50    & $7\!\times\!10^{-4}$ & 2000 & $5\!\times\!10^{-5}$ & 2000 & 512 & $10^{-3}$ cos \\
        & CUB  & 671  & ResNet-18    & $2\!\times\!10^{-4}$ & 4000 & $5\!\times\!10^{-5}$ & 4000 & 512 & $10^{-3}$ cos \\
        & P365 & \chk{2186} & \chk{ResNet-50} & $7\!\times\!10^{-4}$ & 200 & \chk{} & \chk{} & 512 & $10^{-3}$ cos \\
\bottomrule
\end{tabular}
\end{table}

\paragraph{Selecting $\lambda_{\text{gate}}$.}
\chk{State the grid and the criterion, e.g.\ ``the largest $\lambda_{\text{gate}}$ in
$\{\dots\}$ whose held-out accuracy is within $x$ points of the baseline'', and which split it
is selected on. For DN-CBM this is the 20\% held-out split; LF-CBM and UCBM have no separate
validation split in these scripts.}

\paragraph{Software and compute.}
\chk{GPU model, PyTorch/CUDA versions, and wall-clock per run.} There is one conda environment
per pipeline, because the original repositories pin incompatible dependencies. Code, configs,
and the single launcher that reproduces every seed sweep (\texttt{run\_seed\_sweeps.py}) will
be released.
